%% =====================================================================
%% THESIS INFORMATION
%% Fill this in once. It is used on the title pages, the committee page
%% and colophon, the curriculum vitae and the PDF metadata.
%% Check everything here again RIGHT BEFORE sending the thesis to the printer.
%% =====================================================================

%% Title
\newcommand{\thesistitle}{Title of the Dissertation}
\newcommand{\thesissubtitle}{Subtitle}
%% The PDF metadata use "title: subtitle". If the title holds commands (a line break \\, maths), give a
%% plain-text version: \newcommand{\thesispdftitle}{Title of the Dissertation: Subtitle}

%% Author
\newcommand{\thesisauthorfirstnames}{Andrea}             % all first names, as on your passport
\newcommand{\thesisauthorlastname}{Bettini}
\newcommand{\thesisauthorinitials}{A.}
\newcommand{\thesisdegree}{Master of Science in Aerospace Engineering,\\Delft University of Technology, the Netherlands} % highest academic title, university, country
\newcommand{\thesisbirthplace}{[town/city], [country]}

%% Defence
\newcommand{\thesisdefencedate}{[weekday] [day] [month] [year] at [hh:mm] o'clock} % e.g. Friday 4 June 2027 at 10:00 o'clock
\newcommand{\thesisyear}{\the\year}  % the year of the defence (e.g. 2027): copyright line, PDF metadata, citation files; until then the year of compiling
\newcommand{\rectormagnificus}{prof.\,dr.\,ir.\ H.\ Bijl} % Rector Magnificus since 15 January 2026; check before printing

%% Doctoral committee, in the prescribed order (the Rector Magnificus, chairperson, is added at the top):
%%   \member[role]{name}{affiliation}   role: promotor or copromotor (in italics), or none
%%   \membergroup{Independent members}  a heading over the members that follow
%% The promotors are also named in the propositions ("approved by the promotor prof.\,dr.\ S.\ Hickel",
%% in English and Dutch), taken from the members with the role promotor.
\newcommand{\thesiscommittee}{%
    \member[promotor]{Prof.\,dr.\ S.\ Hickel}{TU Delft}
    \member[copromotor]{Dr.\ S.\ Hulshoff}{TU Delft}
    \membergroup{Independent members}
    \member{Prof.\,dr.\ A.\ Jansen}{TU Delft}
    \member{Prof.\,dr.\ N.\ Nescio}{Politecnico di Milano, Italy}
    \member{Prof.\,dr.\,ir.\ J.\ Doe}{TU Delft, reserve member}
    \membergroup{Other members}
    \member{Dr.\ R.\,P.\ Dwight}{TU Delft}
    \member{Dr.\ M.\ Pini}{TU Delft}
}
%% Optional: a sentence about someone who contributed greatly, or a funding statement. Leave empty {} to omit.
\newcommand{\thesiscommitteenote}{}
\newcommand{\thesisdedication}{[For \ldots]}                         % a short line on its own page after the colophon; {} for none
\newcommand{\thesisfunding}{} % e.g. This research was funded by the Dutch Research Council (NWO) under project number ...

%% Colophon
\newcommand{\thesiskeywords}{keyword one, keyword two, keyword three}
\newcommand{\thesisprinter}{[printer]}
%% For the printer (the printer edition, cover.tex and propositions.tex). Ask them:
\newcommand{\thesisprintstandard}{X-6}  % PDF/X-6 (on PDF 2.0), or X-4 if the printer cannot take PDF 2.0 yet
\newcommand{\thesisprintprofile}{}      % the printer's colour profile (an .icc file, e.g. in the main folder); {} for FOGRA39 (coated paper)
\newcommand{\thesisbleed}{2mm}          % bleed around each page of the thesis: often 2 or 3 mm (the cover's is set in settings/cover.tex)
\newcommand{\thesiscover}{[cover designer]}
\newcommand{\thesisisbn}{[978-94-...]}
\newcommand{\thesisdoi}{}  % the DOI the TU Delft repository gives the dissertation (10.4233/uuid:...), once known; {} until then
\newcommand{\thesisorcid}{0000-0000-0000-0000}                  % your ORCID iD (https://orcid.org), in the CV; {} to omit

%% Licence of the thesis. TU Delft publishes theses open access in its repository; a Creative Commons
%% licence says how others may reuse it. CC BY 4.0 is the most common choice. One of:
%%   CC BY 4.0, CC BY-SA 4.0, CC BY-ND 4.0, CC BY-NC 4.0, CC BY-NC-SA 4.0, CC BY-NC-ND 4.0, CC0 1.0
%% or empty {} for "all rights reserved". It sets the sentence on the colophon, its link and the licence in
%% the PDF's metadata. Figures reproduced from published papers may still need the publisher's permission.
\newcommand{\thesislicence}{CC BY 4.0}

%% Data and code availability: where the data and code behind the thesis are kept (TU Delft's research
%% data policy and funders such as NWO expect this), e.g. 4TU.ResearchData or Zenodo, with a DOI. Each
%% gets its own row on the colophon. Cite them in the chapters as well (@dataset and @software in
%% references.bib). Write a DOI as \doi{<doi>}: typewriter on paper, as in the references, and a link on
%% screen. Leave either empty {} to omit its row.
\newcommand{\thesisdata}{4TU.ResearchData, \doi{10.4121/00000000}}
\newcommand{\thesiscode}{Zenodo, \doi{10.5281/zenodo.0000000}}
